% !TEX root = sga4.5-en.tex
% ENGLISH EDITION -- TRANSLATED FROM FRENCH (AI-assisted, needs human review).

\chapter{The cohomology class associated to a cycle}\label{IV}
\blfootnote{by A.\ Grothendieck, written up by P.\ Deligne}




This expos\'e is inspired by notes by Grothendieck, which formed a state 0
of \cite[IV]{sga5}. We define there the cohomology class of a cycle $X$ in
a separated smooth scheme of finite type over a field and we prove that
intersection corresponds to cup-product.

Chapter \ref{IV:1} contains some general sorites. In Chapter \ref{IV:2}, one
defines the class of a cycle, in several situations more general
than the one mentioned above. The main compatibility not considered is
that between direct image of a cycle and trace morphism in cohomology. In
Chapter \ref{IV:3}, one deduces from this formalism the Lefschetz trace formula for
an endomorphism with isolated fixed points of a proper smooth scheme over $k$
algebraically closed -- and for a Frobenius endomorphism of a curve.

We make the following assumptions:
\begin{enumerate}[\indent 1)]
  \item ``scheme'' means noetherian separated scheme (this is
    largely a convenience hypothesis).
  \item In Chapter \ref{IV:2}, one fixes an integer $n$, and $n$ is invertible on
    all the schemes considered.
  \item In Chapter \ref{IV:3}, one fixes a prime $\ell$, and $\ell$ is
    invertible on all the schemes considered. The cohomology used
    is always $\ell$-adic cohomology:
    \[
      \h^\bullet(X) = \dQ_\ell\otimes_{\dZ_\ell} \varprojlim \h^\bullet(X,\dZ/\ell^n) \text{.}
    \]
\end{enumerate}
Cohomology classes of cycles, trace morphisms, \ldots are defined by
passage to the limit from the case of finite coefficients $\dZ/\ell^n$ (cf.
\cite[VI]{sga5}).


















\section{Cohomology with support and cup--products}\label{IV:1}

This paragraph contains reminders of general topology, which the
reader is invited to consult only as needed.










\subsection{\texorpdfstring{$H^1$}{H1} and torsors}\label{IV:1-1}




\subsubsection{}\label{IV:1-1-1}

Let $\sF$ be an abelian sheaf on a site $X$. We know that $\h^1(X,\sF)$
classifies $\sF$-torsors on $X$. We normalize ($=$choose the
sign of) the isomorphism $($set of isomorphism classes of
$F$-torsors$)\to \h^1(X,\sF)$ so that for every exact sequence
$0\to \sF\xrightarrow\alpha\sG\xrightarrow\beta\sH\to 0$ and every
$h\in\h^0(X,\sH)$, the $\sF$-torsor $\beta^{-1}(h)\subset \sG$, on which
$F$ acts by $(f,x)\mapsto \alpha(f)+x$, is of class $\partial h$.




\subsubsection{}\label{IV:1-1-2}

Let $\sP$ be an $\sF$-torsor. If $(U_i)$ is an open covering of $X$, and
$p_i$ a section of $\sP$ over $U_i$, one associates to $\sP$ the
\v{C}ech cocycle
\[
  p_{i j} = p_j - p_i \qquad \text{($p_{i j}\in \h^0(U_i\times U_j,\sF)$).}
\]
If, according to the usual rule, one defines the map
$\operatorname{\check H}^\bullet(X,\sF)\to \h^\bullet(X,\sF)$ so that
it is a morphism of $\delta$-functors, the image of
$(p_{i j})\in \operatorname{\check H}^1(X,\sF)$ in $\h^1(X,\sF)$ is the
class of $\sP$, as normalized in \ref{IV:1-1-1}.




\subsubsection{}\label{IV:1-1-3}

The definition of $\h^\bullet(X,\sF)$ is as follows: for $\sF^\bullet$ a
resolution with acyclic components of $\sF$,
$\h^i(X,\sF) = \h^i\Gamma(X,\sF^\bullet)$. The $\delta$-functor
structure is obtained by associating to a short exact sequence of sheaves a short
exact sequence of resolutions which remains exact after applying the
functor $\Gamma$. If $\sF^\bullet$ is a resolution of $\sF$,
the connecting homomorphism $\partial$ associated to
\[\xymatrix{
  0 \ar[r]
    & \sF \ar[r]
    & \sF^0 \ar[r]^-d
    & \ker(d) \ar[r]
    & 0
}\]
induces the opposite of the defining isomorphism
$\h^1(X,\sF) = \Gamma\left(X,\ker(d)\right)/d \Gamma(X,\sF^0)$.




\subsubsection{}\label{IV:1-1-4}

Let $U$ be an open subset of $X$ (a subsheaf of the final sheaf) and
let $D$ be the ``complementary closed subset.'' We know that $\h_D^1(X,\sF)$
classifies $\sF$-torsors on $X$, trivialized over $U$. For every short
exact sequence $0\to \sF\xrightarrow\alpha \sG\xrightarrow\beta\sH\to 0$, and
every section with support in $D$ $h\in\h_D^0(X,\sH)$, the torsor
$\beta^{-1}(h)$, trivialized over $U$ by the section $0$, has class
$\partial h$.

The long exact sequence of cohomology with support
\[\xymatrix{
  \cdots \ar[r]^-\partial
    & \h_D^i(X,\sF) \ar[r]
    & \h^i(X,\sF) \ar[r]
    & \h^i(U,\sF) \ar[r]^-\partial
    & \cdots
}\]
is defined from the sequence of functors
\[\xymatrix{
  0 \ar[r]
    & \Gamma_D \ar[r]
    & \Gamma \ar[r]
    & \Gamma(U,-) \ar[r]
    & 0
}\]
(exact on injective sheaves). For every section $f\in \h^0(U,\sF)$,
$\partial f\in \h_D^1(X,\sF)$ is the class of the trivial torsor $\sF$,
trivialized over $U$ by the section $f$.




\subsubsection{}\label{IV:1-1-5}

Let $j:U\hookrightarrow X$. If $\sF$ injects into $j_* j^* \sF$, the exact
sequence
\[\xymatrix{
  0 \ar[r]
  & \sF \ar[r]
  & j_* j^* \sF \ar[r]
  & j_* j^* \sF / \sF \ar[r]
  & 0
}\]
yields
$\partial:\h^0(j_* j^*\sF/\sF) = \h_D^0(j_* j^* \sF/\sF) \to \h_D^1(X,\sF)$.
The composite map
\[\xymatrix{
  \h^0(U,\sF) = \h^0(X,j_* j^*\sF) \ar[r]
    & \h^0(X,j_* j^*\sF / \sF) \ar[r]^-\partial
    & \h_D^1(X,\sF)
}\]
is the opposite of the map considered in \ref{IV:1-1-4}.




\subsubsection{}\label{IV:1-1-6}

Recall that if $\sL$ is an invertible sheaf on a scheme $X$, the
$\dG_m$-torsor corresponding to it is the sheaf $\operatorname{Isom}(\sO,\sL)$,
on which $\dG_m$ acts by
$(\lambda,f)\mapsto f\circ (\lambda\cdot) = \lambda f$. Recall also that
if $D$ is a Cartier divisor on $X$, and $j:U\hookrightarrow X$
the complementary open inclusion, the invertible sheaf $\sO(D)$ is the
subsheaf of $j_*\sO_U$ consisting of local sections $s$ such that $s f$
lies in $\sO_X$, for $f$ a local equation of $D$.










\subsection{Cup-products}\label{IV:1-2}

In this number, we develop some remarks on cup-products in
cohomology with support which we will use, in [\nameref{V}], to relate
Poincar\'e duality of curves and self-duality of the Jacobian.

\subsubsection{}\label{IV:1-2-1}

Let $X$ be a site, $Y$ a closed subset of $X$, and $\sF,\sG$ two abelian
sheaves on $X$. For example: $X$ a scheme, $Y$ a closed subscheme
and $\sF,\sG$ sheaves on $\et X$. Let us define a product
\begin{equation*}\tag{1.2.1.1}\label{IV:eq:1-2-1-1}
  \Gamma_Y(X,\sF)\otimes \Gamma(Y,\sG) \to \Gamma_Y(X,\sF\otimes\sG) \text{.}
\end{equation*}
The product of a section $s$, with support in $Y$, of $\sF$ by a section
$t$ of $\sG$ over $Y$ is obtained as follows: locally on $X$, $t$ is the
restriction to $Y$ of a section $t'$ of $\sG$ on $X$, and one forms the product
$s\otimes t'$. It has support in $Y$, and does not depend on the choice of
$t'$, which justifies and allows globalizing the definition.

Let $i:Y\hookrightarrow X$ be the inclusion morphism. The local analogue of
\eqref{IV:eq:1-2-1-1} is the product.
\begin{equation*}\tag{1.2.1.2}\label{IV:eq:1-2-1-2}
  i^!\sF\otimes i^*\sG \to i^!(\sF\otimes\sG) \text{,}
\end{equation*}
from which \eqref{IV:eq:1-2-1-1} is deduced by applying $\Gamma(Y,-)$.




\subsubsection{}\label{IV:1-2-2}

Let us derive these maps. Let $\sK$, $\sL$ and $\sM$ be in the derived
category, and a bilinear map $\sK\lotimes\sL\to\sM$. By
deriving \eqref{IV:eq:1-2-1-1}, one deduces
\begin{equation*}\tag{1.2.2.1}\label{IV:eq:1-2-2-1}
  \eR\Gamma_Y(X,\sK) \lotimes \eR\Gamma(Y,\sL) \to \eR\Gamma_Y(X,\sM)
\end{equation*}
inducing
\begin{equation*}\tag{1.2.2.2}\label{IV:eq:1-2-2-2}
  \h_Y^i(X,\sK)\otimes \h^j(Y,\sL) \to \h_Y^{i+j}(X,\sM)
\end{equation*}
(the cup-product). The local map \eqref{IV:eq:1-2-1-2} yields
\begin{equation*}\tag{1.2.2.3}\label{IV:eq:1-2-2-3}
  \eR i^! \sK\lotimes i^*\sL \to \eR i^!\sM \text{,}
\end{equation*}
from which \eqref{IV:eq:1-2-2-1} is deduced by applying $\eR\Gamma(Y,-)$.




\subsubsection{}\label{IV:1-2-3}

Above, I have glossed over the conflicts entailed by the use in one
same formula of right derivations ($\eR\Gamma$) and left ones
($\lotimes$).
\begin{enumerate}[\indent a)]
  \item In the next number, the right derived functors considered will all
    be of finite cohomological dimension. This allows working
    systematically in the derived categories $\eD^-$. To have
    complexes that are both flat and flasque, one uses the canonical
    flasque resolutions as in \cite[XVII]{sga4}.
  \item For a more general theory, it is no longer possible to interpret
    bilinear maps from $\sK$ and $\sL$ to $\sM$ as
    morphisms from $\sK\lotimes\sL$ to $\sM$. For example,
    $\eR\Gamma_Y(X,\sK)\lotimes\eR\Gamma(Y,\sL)$ is not defined if both
    factors lie in $\eD^+$. One solution is to work in $\eD^+$, to
    define
    \[
      \operatorname{Bil}(\sK,\sL;\sM) = \varinjlim \hom(\sK'\otimes\sL',\sM') \text{,}
    \]
    where the limit is taken over quasi-isomorphisms $\sK'\iso \sK$,
    $\sL'\iso \sL$, $\sM\iso \sM'$ ($\sK'$, $\sL'$, $\sM'$ bounded
    below) and where $\hom$ means ``morphism of complexes, up
    to homotopy,'' and to use systematically such bilinear
    maps, without ever mentioning $\otimes$.
\end{enumerate}




\subsubsection{Second theme}\label{IV:1-2-4}

Let $U$ be an open subset of $X$ and $j:U\hookrightarrow X$ the inclusion morphism.
For $\sK$ in the derived category, on $U$, one puts
$\eR\Gamma_!(U,\sK)=\eR\Gamma(X,j_!\sK)$. For $\sK$, $\sL$ and $\sM$ on $U$, and
a bilinear map $\sK\lotimes\sL\to\sM$, one wants to define
\begin{equation*}\tag{1.2.4.1}\label{IV:eq:1-2-4-1}
  \eR\Gamma(U,\sK)\lotimes\eR\Gamma_!(U,\sL)\to\eR\Gamma_!(U,\sM) \text{.}
\end{equation*}
At the level of sheaves, and of their global sections, such a product is
deduced from the isomorphism
$j_*\sF\otimes j_!\sG\xleftarrow\sim j_!(\sF\otimes \sG)$, but one must beware
of the fact that $\eR\Gamma_!$ is in general not the derived functor of the
functor $\Gamma(X,j_!-)$.

One starts by defining
\[\xymatrix{
  \eR j_* \sK\lotimes j_! \sL
    & \ar[l]_-\sim j_!(\sK\lotimes \sL) \ar[r]
    & j_!\sM \text{.}
}\]
Applying $\eR\Gamma(X,-)$, one finds
\[\xymatrix{
  \eR\Gamma(U,\sK)\lotimes\eR\Gamma_!(U,\sL) = \eR\Gamma(X,\eR j_*\sK)\lotimes \eR\Gamma(X,j_!\sL) \ar[r]
    & \eR\Gamma(X,j_! \sM) \text{.}
}\]

Here again, the conflict between left and right will be resolved in the next
number by working in $\eD^-$.




\subsubsection{Coda}\label{IV:1-2-5}

Let $j:U\hookrightarrow X$ be an open subset and $i:Y\hookrightarrow U$ a
closed subset of $U$. Let $\bar Y$ be a closed subset of $X$ such that
$\bar Y\cap U=Y$, (for example, the complement of $U\setminus Y$). Let $\sK$
be on $U$. One puts $\eR\Gamma_{Y!}(U,\sK)=\eR\Gamma_!(Y,\eR i^!\sK)$ where
$\eR\Gamma_!$ is relative to the inclusion of $Y$ in $\bar Y$. For every open
$V$ of $U$, containing $Y$, $\eR\Gamma_{Y!}(U,\sK)$ maps to
$\eR\Gamma_!(V,\sK)$: if one still denotes by $i$ the inclusion of $\bar Y$ in $X$,
$j$ that of $Y$ in $\bar Y$ and $k$ that of $V$ in $X$, one has
$\eR i^!\sK = j^*\eR i^! k_!(k^*\sK)$, whence a morphism
$j_!\eR i^!\sK \to \eR i^! k_!(k^*\sK)$. Applying
$\eR\Gamma(\bar Y,-)$ to it, one finds
$\eR\Gamma_{Y!}(U,\sK)\to \eR\Gamma_{\bar Y}(k_! k^* \sK) \to \eR\Gamma(k_! k^*\sK) = \eR\Gamma_!(V,\sK)$.

For $\sK,\sL,\sM$ on $U$, and a bilinear map
$\sK\lotimes\sL\to\sM$, one wants to define
\begin{equation*}\tag{1.2.5.1}\label{IV:eq:1-2-5-1}
  \h^n(U,\sK)\otimes \h_!^m(Y,\sL) \to \h_{Y!}^{n+m}(U,\sM)
\end{equation*}
(this last group itself mapping to $\h_!^{n+m}(V,\sM)$). In the
derived category, this amounts to defining
\begin{equation*}\tag{1.2.5.2}\label{IV:eq:1-2-5-2}
  \eR\Gamma_Y(U,\sK) \lotimes \eR\Gamma_!(R,\sL) \to \eR\Gamma_{Y!}(U,\sM) \text{.}
\end{equation*}
One identifies $\eR\Gamma_Y$ with $\eR\Gamma(Y,\eR i^!\sK)$. The desired product
is then of type \eqref{IV:eq:1-2-4-1} relative to the local product
\eqref{IV:eq:1-2-2-3} on $Y$: $\eR i^! \sK\lotimes\sL \to \eR i^! \sM$.




\subsubsection{}\label{IV:1-2-6}

Above, we have unwound the absolute sorites. There is a parallel relative
sorites, with $\Gamma$ replaced by $f_*$ for $f$ a morphism $X\to S$.










\subsection{The Koszul rule}\label{IV:1-3}

Let $A$ be a commutative ring, and $(V_i)_{i\in I}$ a finite family of
graded (or $\dZ/2$-graded) $A$-modules. Recall the definition of the
graded tensor product $\bigotimes_{i\in I} V_i$, in the sense of the
Koszul rule (cf. \cite[XVII 1.1]{sga4}). For each total order $a$ on $I$, we
shall define a module $V(a)$. We shall also define a transitive system
of isomorphisms $\varphi_{a b}:V(b)\iso V(a)$ and the graded
tensor product of the $V_i$ will be the ``common value'' $\varprojlim V(a)$ of this
system of modules. We take
\begin{enumerate}[\indent a)]
  \item $V(a) = \bigotimes_{i\in I} |V_i|$ (ordinary tensor product of the
    ungraded modules underlying the $V_i$).
  \item if the $x_i\in V_i$ are homogeneous, one takes
    $\varphi_{a b}(\otimes x_i) = (-1)^N \otimes x_i$, where $N$ is the sum
    of the $\deg(x_i)\deg(x_j)$ extended over pairs $(i,j)$ such that
    $i<_a j$ and $i>_b j$.
\end{enumerate}




\subsubsection*{Example 1}

Take $I=\{1,2\}$ and let $a$ be the order where $1<2$, $b$ the order where $1>2$.
Let $v_i\in V_i$, homogeneous, and denote by $v_1\otimes v_2$ (resp.
$v_2\otimes v_1$) the image in the graded tensor product of the product of the
$v_i$ in $V(a)$ (resp. $V(b)$). One has
\[
  v_1\otimes v_2 = (-1)^{\deg(v_1)\deg(v_2)} v_2\otimes v_1
\]
(Koszul rule).




\subsubsection*{Example 2}

If the $V_i$ are all of degree $1$, one has, relating the underlying module of the
graded tensor product, and the ordinary tensor product of the
$|V_i|$ underlying the $V_i$, a canonical isomorphism
\[
  \left|\bigotimes V_i\right| \simeq \bigotimes |V_i|\otimes_{\dZ}\bigwedge^{|I|}\dZ^I \text{.}
\]




\subsubsection*{Example 3}

For $A$ a field, and $X_i$ a finite family of spaces, the
K\"unneth formula reads
$\h^\bullet\left(\prod X_i,A\right) = \bigotimes \h^\bullet(X_i,A)$.

Let $V^\vee$ be the graded dual of $V$. The canonical map
\begin{equation*}\tag{1.3.1}\label{IV:eq:1-3-1}
  V^\vee\otimes V = V\otimes V^\vee \to A
\end{equation*}
is $v'\otimes v\to v'(v)$.

Assume now that $A$ is a field, and consider only
finite-dimensional vector spaces. The canonical isomorphism
\begin{equation*}\tag{1.3.2}\label{IV:eq:1-3-2}
  W\otimes V^\vee \to \hom(V,W)
\end{equation*}
is $w\otimes v'\mapsto \left(v\mapsto w\cdot v'(v)\right)$. Via this
isomorphism, the composition $\hom(Y,Z)\otimes \hom(X,Y)\to \hom(X,Z)$
is identified with the morphism induced by \eqref{IV:eq:1-3-1}:
$Z\otimes Y^\vee\otimes Y\otimes X^\vee \to Z\otimes X^\vee$.

The trace of $f:V\to V^\vee$ (zero for $f$ homogeneous of degree $\ne 0$) is
the image of $f$ under
\begin{equation*}\tag{1.3.4}\label{IV:eq:1-3-4}
  \tr:\hom(V,V) {\xleftarrow\sim}_\text{\eqref{IV:eq:1-3-2}} V\otimes V^\vee = V^\vee \otimes V \to_\text{\eqref{IV:eq:1-3-1}} A .
\end{equation*}

One checks easily that, for $f$ of degree $0$,
\begin{equation*}\tag{1.3.5}\label{IV:eq:1-3-5}
  \tr(f,V) = \sum (-1)^i \tr(f,V^i) \text{.}
\end{equation*}

Expressing that the two composite morphisms of morphisms
\eqref{IV:eq:1-3-1} $V^\vee\otimes V\otimes W^\vee\otimes W\to k$ commute, one
finds that, for $f:V\to W$ and $g:W\to V$ homogeneous, one has
\begin{equation*}\tag{1.3.6}\label{IV:eq:1-3-6}
  \tr(f g) = (-1)^{\deg(f)\deg(g)} \tr(g f) \text{.}
\end{equation*}


















\section{The cohomology class associated to a cycle}\label{IV:2}









\subsection{The class of a divisor}\label{IV:2-1}




\subsubsection{}\label{IV:2-1-1}

Let $D$ be a Cartier divisor in a scheme $X$. Outside $D$, the invertible
sheaf $\sO(D)$ is trivialized by the section $1$. The class
$\cl(D)$ of $D$, in $\h_D^1(X,\dG_m)$, is the class of the
$\dG_m$-torsor trivialized over $X\setminus D$ corresponding to it
(\ref{IV:1-1-6} and \ref{IV:1-1-4}).

Let $\partial:\h^i(X\setminus D,\dG_m)\to \h_D^i(X,\dG_m)$ be the morphism
\ref{IV:1-1-4}. If $D$ admits a global equation $f$, multiplication by
$f$ is an isomorphism from $\sO(D)$, trivialized by $1$ over $X\setminus D$,
to $\sO$, trivialized by $f$ over $X\setminus D$. By \ref{IV:1-1-4},
one therefore has
\begin{equation*}\tag{2.1.1}\label{IV:eq:2-1-1}
  \cl(D) = \partial f \text{.}
\end{equation*}

For every morphism $u:X'\to X$ such that $u^* D$ is still a Cartier divisor
(i.e., $u^{-1} D$ disjoint from $\operatorname{Ass}(X')$), one has
$\cl(u^* D) = u^*\cl(D)$. If one wanted such a
functoriality for every morphism $u$, one would have to consider not
Cartier divisors, but more generally invertible sheaves
equipped with a section.

Recall that the integer $n$ is henceforth assumed invertible on the
schemes considered. Let
$\partial:\h_D^i(X,\dG_m) \to \h_D^{i+1}(X,\dmu_n)$ be the boundary map for the
Kummer exact sequence $0 \to \dmu_n \to \dG_m \to \dG_m \to 0$.




\begin{definition}\label{IV:2-1-2}
The \emph{class $\cl_n(D)$} of $D$ in $\h_D^2(X,\dmu_n)$ is
$\partial \cl(D)$.
\end{definition}

When there is no risk of confusion, we shall omit mention of $n$.




\subsubsection{}\label{IV:2-1-3}

The diagram
\[\xymatrix{
  \h^0(X\setminus D,\dG_m) \ar[r]^-\partial \ar[d]^-\partial
    & \h_D^1(X,\dG_m) \ar[d]^-\partial \\
  \h^1(X\setminus D,\dmu_n) \ar[r]^-\partial
    & \h_D^2(X,\dmu_n)
}\]
is anticommutative. If $D$ admits a global equation $f$,
$\cl_n(D)$ is therefore the opposite of the image under $\partial$ of the
class in $\h^1(X\setminus D,\dmu_n)$ of the $\dmu_n$-torsor of $n$-th
roots of $f$.




\begin{proposition}\label{IV:2-1-4}
Let $i$ be the inclusion of $D$ in $X$. If $D$ and $X$ are regular, the
cohomology sheaves with support $\eR^p i^! \dmu_n$ vanish for $p=0,2$,
and $\eR^2 i^!\dmu_n = \underline{\dZ/n}$, generated by
$\cl_n(D)$.
\end{proposition}

It suffices to prove that, for $X$ strictly local and $D$ defined by a
regular parameter, one has $\h_D^p(X,d\mu_n) = 0$ for $p=0,1$ and
$\h^2(X,\dmu_n) = \dZ/n$ generated by $\cl(D)$. Denoting
reduced cohomology by $\sim$, one has
$\widetilde\h^{p-1}(X\setminus D,\dmu_n)\iso \h_D^p(X,\dmu_n)$. The assertion
for $p=0,1$ expresses that $D$ does not disconnect $X$, and for $p=2$ follows,
via \ref{IV:2-1-3}, from Abhyankar's lemma.

This is a partial analogue of the relative theorem
(\hyperref[I]{Cohomologie \'etale}, \ref{I:5-3-4}). Grothendieck conjectures
that the $\eR^p i^! \dmu_n$ vanish for $p\ne 2$, at least for $X$ excellent
(purity conjecture), but this is known only in characteristic $0$
(\cite[XIX]{sga4}).




\begin{theorem}[Fundamental compatibility]\label{IV:2-1-5}
Let $X$ be a smooth curve over an algebraically closed field $k$, $P$ a
closed point of $X$ and $\tr$ the composite $\h_P^2(X,\dmu_n) \to \h_c^2(X,\dmu_n) \xrightarrow{\tr} \dZ/n$. One has
\[
  \tr \cl(P) = 1 \text{.}
\]
\end{theorem}

Let $\bar X$ be the smooth projective curve completing $X$. The formula
expresses that the invertible sheaf $\sO(P)$ on $\bar X$ is of degree $1$.










\subsection{Cohomological method}\label{IV:2-2}




\subsubsection{}\label{IV:2-2-1}

Let $X$ be a (noetherian) scheme. Recall that a subscheme $Y$ of $X$
is said to be a local complete intersection, of codimension $c$, if, locally
(on $Y$), it is defined by a regular sequence of $c$ equations in
$X$. For $X$ the spectrum of a local ring $A$ with maximal ideal $\fm$ and
$Y$ with ideal $\fa$, this means that $\exp^i(A/\fa,A) = 0$ for
$i<\dim(\fa/\fm\fa)=c$, and every sequence of elements of $\fa$ whose image in
$\fa/\fm\fa$ is a basis of $\fa/\fm\fa$, is a regular sequence
of equations for $Y$.




\subsubsection{}\label{IV:2-2-2}

Let $i:Y\hookrightarrow X$ be a local complete intersection, of codimension
$c$. We propose to define a local fundamental class
$\cl(Y)$ which is a global section of the cohomology sheaf with support
$\eR^{2c}i^! \dZ/n(c)$ (recall that $\dZ/n(c)=\dmu_n^{\otimes c}$).

Locally, $Y$ is the intersection of a sequence of $c$ divisors $D_i$, and one
defines $\cl(Y)$ as the cup-product of the $\cl(D_i)$. Each $\cl(D_i)$ has
support in $D_i$, their product has support in $Y$. That, locally,
this product does not depend on the choice of the $D_i$, follows from \ref{IV:2-2-3}
below and the following invariance properties:
\begin{enumerate}[\indent a)]
  \item compatibility with localization;
  \item independence of the order of the $D_i$ (the $\cl(D_i)$ are of even
    degree $2$, so the cup-product is commutative);
  \item the product depends only on the ``flag''
    $D_1\supset D_1\cap D_2\supset \cdots\supset Y$.
\end{enumerate}
To prove c), one notes the existence of a product (variant of \ref{IV:1-2-1})
\[
  \h_{D_1}^\bullet(X)\otimes \h_{D_1\cap D_2}^\bullet(D_1)\otimes \cdots \otimes \h_Y^\bullet(D_1\cap \cdots \cap D_{c-1}) \to \h_Y^\bullet(X) \text{;}
\]
the product of the $\cl(D_i)$ is still the product of the ($\cl(D_i)$ restricted to
$D_1\cap \cdots \cap D_{i-1})=(\cl(D_1\cap \cdots \cap D_i)$ in
$D_1\cap \cdots \cap D_{i-1})$.




\begin{lemma}\label{IV:2-2-3}
Let $A$ be a local ring with maximal ideal $\fm$ and $u=(u_1,\dots,u_c)$,
$v=(v_1,\dots,v_c)$ two regular sequences generating the same ideal
$\fa$. There then exists a sequence $w_i$ ($1\leqslant i\leqslant N$) of such
sequences, joining them, such that $w_{i+1}$ is deduced from $w_i$ by a
permutation, or by changing only the last element.
\end{lemma}

Since permutations are allowed, it would amount to the same to allow
changing a single element, rather than the last one. By \ref{IV:2-2-1}, one is
then reduced to checking that, in the vector space $\fa/\fm\fa$, one
can pass from one basis to another by a sequence of permutations and of
base changes modifying only one vector. The linear group is
indeed generated by diagonal, elementary and
permutation matrices.




\subsubsection{}\label{IV:2-2-4}

The same methods allow defining a local fundamental class
for every $Y\subset X$ locally definable by $c$ equations (where
fewer equations suffice, the class is zero).




\subsubsection{}\label{IV:2-2-5}

One can pass from such a local fundamental class, in
$\h^0\left(Y,\eR^{2 c} i^!\dZ/n(c)\right)$, to a global fundamental class,
in $\h_Y^{2 c}\left(X,\dZ/n(c)\right)$, when semi-purity results are
available:




\begin{proposition}\label{IV:2-2-6}
\begin{enumerate}[(i)]
  \item If $\eR^p i^!\dZ/n=0$ for $p<2 c$, then
    \[
      \h_Y^{2 c}\left(X,\dZ/n(c)\right) \iso \h^0\left(Y,\eR^{2 c} i^! \dZ/n(c)\right) \text{.}
    \]
  \item Let $Z$ be a closed subset of $Y$, with complement $V$ in $Y$, and
    $k$ the inclusion of $Z$ in $X$. If $\eR^p i^! \dZ/n = 0$ for
    $p\leqslant 2 c$, then
    \[
      \h_Y^{2 c}\left(X,\dZ/n(c)\right) \hookrightarrow \h_V^{2 c}\left(X,\dZ/n(c)\right) \text{.}
    \]
    If $\eR^p i^!\dZ/n=0$ for $p\leqslant 2 c+1$, this map is an
    isomorphism.
\end{enumerate}
\end{proposition}

The hypothesis $\eR^p i^!\dZ/n=0$ is equivalent to $\eR^p i^!\dZ/n(c)=0$. This said,
(i) is read on the spectral sequence
$\h^p(Y,\eR^q i^!)\Rightarrow \h_Y^{p+q}(X,-)$. If $k_1$ denotes the inclusion of $Z$ in
$Y$, one has $k=i k_1$, whence $\eR k^! \dZ/n=\eR k_1^! \eR i^!$, and the long exact
cohomology sequence for $Z\subset Y$ yields a long exact sequence
\[\xymatrix{
  \cdots \ar[r]
    & \h^i\left(Z,\eR k^! \dZ/n(c)\right) \ar[r]
    & \h_Y^i\left(X,\dZ/n(c)\right) \ar[r]
    & \h_V^i\left(X,\dZ/n(c)\right) \ar[r]
    & \cdots
}\]
from which (ii) follows.




\subsubsection{Amplification}\label{IV:2-2-7}

Proposition \ref{IV:2-2-6} remains valid for $i$ any
separated morphism of finite type, and $2 c$ any integer (positive or negative),
provided one replaces $\h_Y^{2 c}\left(X,\dZ/n(c)\right)$ by
$\h^{2 c}\left(Y,\eR i^!\dZ/n(c)\right)$, and likewise for $\h_V$.

The following semi-purity results follow from
\cite[1.8, 1.10, 1.15]{sga2}. We shall recall their proof.




\begin{theorem}\label{IV:2-2-8}
Let
\[\xymatrix{
  X \ar[r]^-i \ar[dr]_-f
    & X \ar[d]^-g \\
  & S
}\]
be a morphism of $S$-schemes of finite type. Assume $X$ smooth, purely of relative dimension $N$ and $Y$ fiberwise
of dimension $\leqslant d$. Let $c=N-d$.
\begin{enumerate}[\indent (i)]
  \item One has $\eR^p i^!\dZ/n=0$ for $p<2 c$; likewise with $\dZ/n$ replaced
    by a sheaf $g^* \sF$.
  \item If, over a dense open $U$ of $S$, $Y$ is fiberwise of
    dimension $<d$, one has $\eR^p i^!\dZ/n=0$ for $p\leqslant 2 c$. If moreover the
    complement $Z$ of $U$ does not locally disconnect $S$, one has
    $\eR^pi^!\dZ/n=0$ for $p\leqslant 2 c+1$.
\end{enumerate}
\end{theorem}

Since $\eR g^! \sF=g^*\sF(N)[2N]$, the transitivity formula
$\eR i^!\eR g^!=\eR f^!$ shows that
$\eR^{2 c+q} i^!(g^*\sF)=0\Leftrightarrow \eR^{-2 d+q}f^!(\sF)$. This
reduces us to studying $f$, i.e. to assuming $X=S$. Assertion (i) is
then \cite[XVIII 3.17]{sga4}.

Let $Y'$ be the inverse image of $Z$ in $Y$:
\[\xymatrix{
  Y' \ar[r]^-v \ar[d]^-f
    & Y \ar[d]^-f \\
  Z \ar[r]^-u
    & S
}\]
By (i), the $\eR^p f^! \dZ/n$ are supported in $Y'$ for
$p\leqslant -2 d+1$; and the spectral sequence
$\eR^a v^!\eR^b f^! \Rightarrow \eR^{a+b} (f v)^!$ shows that they coincide with
the $\eR^p(f v)^!\dZ/n=\eR^p(u f)^!\dZ/n$. Applying (i) to $Y'/Z$ and to the
spectral sequence $\eR^a f^!\eR^b u^! \Rightarrow \eR^{a+b} (u f)^!$, one finds that
(ii) follows from the vanishing of $\eR^b u^! \dZ/n$ for $b=0$, or $b=0$ and $1$
according to the case.




\subsubsection{}\label{IV:2-2-9}

Grothendieck conjectures the following absolute analogue of \ref{IV:2-2-8} (semi-purity
conjecture -- a consequence of the purity conjecture): for $Y$
of codimension $\geqslant c$ in $X$ regular, one has $\h_Y^i(X,\dZ/n)=0$ for
$i<2 c$, at least if $X$ is excellent.




\subsubsection{}\label{IV:2-2-10}

We are now ready to define the class of a cycle $Y$ of
codimension $c$ in $X$ smooth over a field $k$. One writes $Y=\sum d_i Y_i$,
where the $Y_i$ are reduced irreducible. An open $U_i$ of $Y_i$, with
complement of codimension $>c$, is then a local complete
intersection in $X$. This allows defining the local fundamental class of
$U_i$. By \ref{IV:2-2-6} and \ref{IV:2-2-8}, it comes from a
unique fundamental class $\cl(Y_i)\in \h_{Y_i}^{2 c}\left(X,\dZ/n(c)\right)$,
and one puts
\[
  \cl(Y) = \sum d_i \cl(Y_i) \in \h_{|Y|}^{2 c}\left(X,\dZ/n(c)\right) \text{.}
\]




\subsubsection{}\label{IV:2-2-11}

In complex analytic geometry, a construction of Baum, Fulton and
MacPherson \cite{bfm75} allows defining without restriction the
cohomology class of a local complete intersection $Y\subset X$. Suppose $Y$
purely of codimension $c$, and let $\sN$ be the vector bundle on $Y$, the
normal bundle of $Y$ in $X$. Its local sections are those of $(\sI/\sI^2)^\vee$,
where $\sI$ is the ideal sheaf of $Y$. Recall that the sheaf of
$C^\infty$ functions on $X$ is defined locally, in terms of local
embeddings of $X$ in $\dC^n$, as the restriction to $X$ of the quotient of the
sheaf of $C^\infty$ functions on $\dC^n$ by the ideal generated by the
real and imaginary parts of the equations defining $X$. The
bundle $\sN$ extends to a complex $C^\infty$ vector bundle $N$ on a
neighborhood $U$ of $Y$ in $X$, and for $U$ small enough, there exist sections
$f$ of $N$, with zero locus $Y$, and such that, on $Y$, $d f:\sN\to N$
is the identity. Two choices of $N$ and $f$ are homotopic on $Y$ small enough.

The cohomology class $\cl(U)$ of the $0$ section of $N$ (denoted $z:U\to N$)
is defined: $U$ is a local complete intersection in $N$, and
$\eR^i z^!\dZ = 0$ for $i\ne 2 c$, $\eR^{2 c} z^!\dZ = \dZ$. One has
$f^*:\h_U^\bullet(N,\dZ) \to \h_Y^\bullet(X,\dZ)$ and one puts
\[
  \cl(Y) = f^*\cl(U) \text{.}
\]

The same construction works as soon as one has on an analytic subspace $Y$
of $X$ the following normal structure: a locally free sheaf $\sC$ of
rank $c$ on $Y_\text{red}$, and an epimorphism
$\sC\to \sI/\sI\cdot \sI_\text{red}$.










\subsection{Homological method}\label{IV:2-3}




\subsubsection{}\label{IV:2-3-1}

Let $f:Y\to X$ be a flat morphism of finite type, with fibers of dimension
$\leqslant d$. In \cite[XVIII 2.9]{sga4}, we have defined a trace
morphism
\begin{equation*}\tag{2.3.1.1}\label{IV:eq:2-3-1-1}
  \tr_f:\eR^{2 d} f_! \dZ/n(d) \to \dZ/n \text{.}
\end{equation*}
One has $\eR^i f_!\dZ/n(d)=0$ for $i>2 d$, and $\eR^i f^!\dZ/n=0$ for $i<-2 d$.
This, and adjunction between $\eR f_!$ and $\eR f^!$, furnish isomorphisms
\begin{equation*}\tag{2.1.3.2}\label{IV:eq:2-1-3-2}
\begin{aligned}
  \hom\left(\eR^{2 d}f_!\dZ/n(d),\dZ/n\right)
    &= \hom\left(\eR f_!\dZ/n(d),\dZ/n[-2 d]\right) \\
    &= \hom\left(\dZ/n,\eR f^!\dZ/n(-d)[-2 d]\right) \\
    &= \h^0\left(Y,\eR^{-2 d}f^!\dZ/n\right) \text{.}
\end{aligned}
\end{equation*}
At least over $S$ a point, the image of $\tr_f$ in the last two groups
deserves the name of fundamental class of $Y$ (in homology).

We still denote by $\tr_f$ the image of $\tr_f$ in the second group, and the
morphisms deduced from it by functoriality. For example, for $Y/S$
proper, the morphism
\begin{equation*}\tag{2.3.1.3}\label{IV:eq:2-3-1-3}
\xymatrix{
  \h^i\left(Y,\dZ/n(d)\right) = \h^i\left(S,\eR f_!\dZ/n(d)\right) \ar[r]
    & \h^{i-2 d}(S,\dZ/n) \text{.}
}
\end{equation*}

Suppose $Y$ contained in $X$ smooth over $S$, purely of relative
dimension $N$, and put $c=N-d$.
\[\xymatrix{
  Y \ar@{^{(}->}[r]^-i \ar[dr]_-f
    & X \ar[d] \\
  & S
}\]
One has $\eR f^! = \eR i^!\eR g^!$, and $\eR g^! \dZ/n=\dZ/n(-N)[-2 N]$. The
fundamental class of $Y$ is identified with an element of
$\h^0\left(Y,\eR i^!\dZ/n(c)[2 c]\right)=\h_Y^{2 c}\left(X,\dZ/n(c)\right)$, the
class $\cl(Y)$ of $Y$ in $X$. We shall see below that it depends only
on $Y\subset X$, not on the projection of $X$ onto $S$, and that for $Y$
a local complete intersection, it induces the local class of the preceding
number.

Let us make explicit the isomorphism
\[
  \h_Y^{2 c}\left(X,\dZ/n(c)\right) \iso \hom\left(\eR^{2 d}f_!\dZ/n(d),\dZ/n\right)
\]
which turns $\cl(Y)$ into $\tr_f$: via the isomorphisms
\begin{align*}
  \h_Y^{2 c}\left(X,\dZ/n(c)\right)
    &= \h_Y^{2 c}\left(\eR i^!\dZ/n(c)\right) \\
    &= \hom_Y\left(\dZ/n(d)[2 d],\eR i^!\dZ/n(N)[2 N]\right) \\
    &= \hom_X\left(\dZ/n(d)[2 d]_Y,\dZ/n(N)[2 n]\right) \text{,}
\end{align*}
it is the top row of
\[\xymatrix{
  \h_Y^{2 c} \ar[r] \ar[dr]_-{(3)}
    & \hom\left(\eR f_!\dZ/n(d)[2 d],\eR f_! \eR i^!\dZ/n(N)[2 n]\right) \ar[r]^-{(1)} \ar[d]^-{\tr_i} \ar@{}[dr]|-{(2)}
    & \hom\left(\eR f_!\dZ/n(d)[2d],\dZ/n\right) \ar@{=}[d] \\
  & \hom\left(\eR g_!\dZ/n(d)[2 d]_Y,\eR g_!\dZ/n(N)[2 N]\right) \ar[r]^-{\tr_g}
    & \hom\left(\eR g_!\dZ/n(d)[2 d]_Y,\dZ/n\right) \text{,}
}\]
where (1) is the adjunction map $\eR f_!\eR f^!\to \text{id}$. The
commutativity (2) expresses that the isomorphism $\eR f^!=\eR i^!\eR g^!$ is
defined by adjunction. Finally, the map deduced from (3)
\[
  \h_Y^{2c}\left(X,\dZ/n(c)\right) \otimes \eR^{2 d} g_!\dZ/n(d) \to \eR^{2 N} g_!\dZ/n(N)
\]
is interpreted as a cup-product (cf. \ref{IV:1-2}).




\begin{definition}\label{IV:2-3-2}
The class $\cl(Y)\in\h_Y^{2 c}\left(X,\dZ/n(c)\right)$ has for characteristic
property that, for every local section $u$ of $\eR^{2 d}f_! \dZ/n(d)$,
one has
\[
  \tr_f(u) = \tr_g\left(\cl(Y)\smallsmile u\right) \text{.}
\]
\end{definition}

Since $\cl(Y)$ already has an image $\tr_f$ in
$\hom\left(\eR f_!\dZ/n(d)[2 d],\dZ/n\right)$, formula \ref{IV:2-3-2} holds
for the maps deduced from $\tr_f$ by functoriality. For example,
for $X$ proper over $S$ and $u\in \h^\bullet(Y,\dZ/n)$, the formula holds in
$\h^\bullet\left(S,\dZ/n(-d)\right)$. For $v\in \h^\bullet(X,\dZ/n)$, it
gives
\[
  \tr_f(i^* v) = \tr_g\left(\cl(Y)\smallsmile v\right) \text{,}
\]
where the cup-product can be computed in $\h^\bullet(X,\dZ/n)$.




\subsubsection{}\label{IV:2-3-3}

We shall define the trace morphisms, and hence the class $\cl(Y)$, under more
general hypotheses. Let thus be a diagram
\[\xymatrix{
  Y \ar@{^{(}->}[r] \ar[dr]_-f
    & X \ar[d]^-g \\
  & S
}\]
with $g$ smooth, purely of relative dimension $N$ and $Y$ closed in $X$,
fiberwise of dimension $\leqslant d$. We equip moreover $Y$ with a
``weight'' $\sK$ of the following type: $\sK\in\eD_{\text{parf},X}^b$ is a
bounded complex of sheaves of $\sO$-modules on $X$, of finite $\tor$-dimension
over $S$ (or over $X$, this amounts to the same) and whose cohomology
sheaves are coherent and supported in $Y$. We propose to
define a morphism $\tr_{f,\sK}:\eR^{2 d}f_! \dZ/n(d) \to \dZ/n$. For $Y$
flat over $S$ and $\sK=\sO_Y$, this will be the preceding trace morphism. In
general, it depends only on the lengths of $\sK$ at the
generic points $y$ of the irreducible components of
$Y(\operatorname{lg}_y(\sK) = \sum (-1)^i \operatorname{lg} \h^i(\sK)_y)$.
% NOTE: notation is not clear, is it a calligraphic H?
Finally, we shall denote by $\cl(Y,\sK)$ the class with support in $Y$ such that
\begin{equation*}\tag{2.3.3.1}\label{IV:eq:2-3-3-1}
  \tr_g\left(\cl(Y,\sK)\smallsmile u\right) = \tr_{f,\sK}(u) \text{.}
\end{equation*}




\subsubsection{}\label{IV:2-3-4}

The construction of $\tr_{f,\sK}$ is parallel to that of
\cite[XVIII.2]{sga4}; we shall only indicate its broad lines.

\paragraph{A}
$d=0$ ($f$ quasi-finite). Let $x$ be a geometric point of $Y$. $s$ its image
in $S$, $\sK_{(x)}$ the inverse image of $\sK$ on the strict localization
of $X$ at $x$, $\sK_{(x)s} = \sK_{(x)}\lotimes_{\sO_{S,s}^h} k(s)$ its inverse
image on the geometric fiber at $s$, and
$n(x)=\sum (-1)^i \dim_{k(s)} \h^i(\sK_{(x)s})$. The function $x\mapsto n(x)$
is a weighting of $f$, and one takes the corresponding trace morphism
\cite[XVII.6.2.5]{sga4}. The weighting $n(x)$, and $\tr_{f,\sK}$ are
compatible with arbitrary base change.

\paragraph{B}
General case. If $u:Z\to \dA_S^d$ is such that $u i$ is quasi-finite, one puts
$\tr_{f,\sK} = \tr_{\dA_S^d}\circ \tr_{u i,\sK}$. This trace morphism is
clearly compatible with any base change $S'/S$. To prove that it does not
depend on $u$, one may therefore assume that $S$ is the spectrum of an
algebraically closed field $k$. Let in this case $Y_i$ be the irreducible
components of $Y$, and $\operatorname{lg}_i$ the length of $\sK$ at the
generic point of $Y_i$. If $\alpha_i$ is the inclusion of
$(Y_i)_\text{red}$ in $Y$, one has
\begin{equation*}\tag{2.3.4}\label{IV:eq:2-3-4}
  \tr_{f,\sK} = \sum \operatorname{lg}_i \tr_{Y_{i,\text{red}/S}} \alpha_i^\ast
\end{equation*}
(this formula follows from the analogous and easy formula for $\tr_{u i}$).

This defines $\tr_{f,\sK}$ locally on $Y$; one then proceeds as in
\cite[XVIII.2.9]{sga4}.




\begin{lemma}\label{IV:2-3-5}
If $g=g'g'':X\xrightarrow{g''} S'\xrightarrow{g'} S$, with $g'$ and $g''$ smooth
and purely of relative dimension $N'$ and $N''$, and $Y$ is still, over
$S'$, fiberwise of codimension $\geqslant c$, then the class
$\operatorname{cl}(Y,\sK)$ is the same, computed in terms of $g$ or of
$g''$.
\end{lemma}

Let $f'=g'' i$ and $d'=d-N'$. The formula $\tr_g=\tr_{g'}\tr_{g''}$
(\cite[XVIII.2.9]{sga4} Var 3) ensures the commutativity of
\[\xymatrix{
  \h_X^{2 c}\left(Z,\dZ/n(c)\right) \ar[d] \\
  \hom_{S'}\left(\eR f_!\dZ/n(d')[2 d'],\eR g_!''\dZ/n(N'')[2 N'']\right) \ar[r] \ar[d]^-{\eR g_!'(N')}
    & \hom_{S'}\left(\eR^{2 d'} f_!' \dZ/n,\dZ/n\right) \ar[d]^-{\tr_{g'}\circ \eR^{2 N'} g_!'} \\
  \hom_S\left(\eR f_!\dZ/n(d)[2d],\eR g_!'\dZ/n(N)[2 N]\right) \ar[r]
    & \hom_S\left(\eR^{2 d} f_!\dZ/n,\dZ/n\right)
}\]
while the analogous and easy-to-check formula
$\tr_{f,\sK}=\tr_{g'}\tr_{f',\sK}$ ensures that the image of $\tr_{f',\sK}$ is
$\tr_{f,\sK}$. The assertion follows.




\begin{lemma}\label{IV:2-3-6}
If $Y$ is a divisor in $X$, the class \ref{IV:2-3-2} coincides with the
class \ref{IV:2-1-2}.
\end{lemma}

By semi-purity (\ref{IV:2-2-6}(i) and \ref{IV:2-2-8}(i)), the problem is
local on $Y$; using \ref{IV:2-3-5} to replace $S$ by a suitable $S'$
this allows us to assume $N=1$. Using semi-purity again
(\ref{IV:2-2-6}(ii) and \ref{IV:2-2-8}(ii)) it suffices to prove
\ref{IV:2-3-6} over the generic points of $S$. An \'etale localization
then reduces us to assuming $S$ the spectrum of an
algebraically closed field $k$. The classes \ref{IV:2-3-2} and \ref{IV:2-1-2} each
being additive in $Y$, this reduces us to the case where $Y$ is a closed
point on a smooth curve over $k$. That the characteristic property
\ref{IV:2-3-2} holds is then the fundamental compatibility
\ref{IV:2-1-5}.




\begin{lemma}\label{IV:2-3-7}
The formation of $\cl(Y,\sK)$ is compatible with any base change
$S'/S$.
\end{lemma}

Follows from the same assertion for the trace morphisms.




\begin{theorem}\label{IV:2-3-8}
(i) $\cl(Y,\sK)$ depends only on $Y\subset X$ and on the lengths of $\sK$ at the
generic points of $Y$. This class is additive in $\sK$. For
$\sK=\sO_Y$ and $Y$ a local complete intersection, it induces the local
class \ref{IV:2-2-2}.

(ii) Let a commutative diagram
\[\xymatrix{
  X' \ar[r]^-u \ar[d]
    & X \ar[d] \\
  S' \ar[r]
    & S
}\]
with $X'$ smooth over $S'$. If $u^{-1}(Y)$ is still fiberwise of
codimension $\geqslant c$, then
\[
  \cl\left(u^{-1}(Y),\eL u^\ast\sK\right) = u^\ast \cl(Y,\sK) \text{.}
\]

(iii) Let $Y'\subset X$ be fiberwise of codimension $\geqslant c'$,
$\sK'$ a weight on $Y'$ and assume that $Y\cap Y'$ is fiberwise of
codimension $\geqslant c+c'$. Then
\[
\cl\left(Y\cap Y',\sK\lotimes \sK'\right) = \cl(Y,\sK) \smallsmile \cl(Y',\sK') \text{.}
\]
\end{theorem}


% The following paragraphs were originally numbered

% (2.3.8.1)
\paragraph{(A)}
Given $Y\subset X/S$ as above, it follows from semi-purity
(\ref{IV:2-2-6}, \ref{IV:2-2-8}) that to check that two classes in
$\h_Y^{2 c}\left(X,\dZ/n(c)\right)$ coincide, it suffices to check it
locally at the generic points of $Y$, or even after any
base change $s\to S$ ($s$ generic geometric point of $S$),
at the generic points of $Y$.

% (2.3.8.2)
\paragraph{(B) Proof of (ii) for $S=S'$ and $u$ a closed immersion}
By localization (A) one may assume that $X'$ is the inverse
image of the $0$ section under a smooth morphism $v:X\to \dA_S^{N'}$
\[\xymatrix{
  X' \ar[r]^-u \ar[d]
    & X \ar[d]^-v \\
  S \ar[r]^-0
    & \dA_S^{N'} \ar[r]^-\pi
    & S
}\]
One applies \ref{IV:2-3-5} to $g=\pi v$ and \ref{IV:2-3-7} to the base
change $S\to \dA_S^{N'}$.

% (2.3.8.3)
\paragraph{(C) In a product situation:}
$X=X'\times_S X''$, $Y=Y'\times_S Y''$, $\sK=\sK\lotimes_S \sK''$, one has
$\cl(Y,\sK) = \cl(Y,\sK')\smallsmile \cl(Y'',\sK'')$ (exterior cup-product,
i.e. $\operatorname{pr}_1^\ast \cl(Y',\sK')\smallsmile \operatorname{pr}_2^\ast \cl(Y'',\sK'')$).

From the functorial properties of $\tr_g$ \cite[XVIII.2.12]{sga4} follows
the commutativity of
\small
\[\xymatrix@=0.5cm{
  \h_{X'}^{2c'}\otimes \h_{X''}^{2 c''}\ar[r] \ar[d]
    & \hom_S\left(\eR f_!' \dZ/n(d')[2 d'],\eR g_!'\dZ/n(N')[2 N']\right)\otimes \cdots \ar[r] \ar[d]
    & \hom_S\left(\eR^{2 d'} f_!' \dZ/n,\dZ/n\right) \otimes \cdots \ar[d] \\
  \h_X^{2 c} \ar[r]
    & \hom_S\left(\eR f_!\dZ/n(d)[2 d],\eR g_!\dZ/n(N)[2 N]\right) \ar[r]
    & \hom_S\left(\eR^{2 d}f_!\dZ/n,\dZ/n\right)
}\]
\normalsize
and one checks that $\tr_f$ is the image of $\tr_{f'}\otimes \tr_{f''}$ whence
the assertion.

\paragraph{Proof of (ii)}
(\ref{IV:2-3-7}) and a preliminary base change reduce us to
assuming $S=S'$. Factor $u$ through its graph:
\[\xymatrix{
  X' \ar[r]^-{(\text{id},u)}
    & X'\times X \ar[r]^-{\operatorname{pr}_2}
    & X \text{.}
}\]
This leads us to treat separately $\operatorname{pr}_2$, which falls
under (C) for $X'=X'$, $\sK'=\sO_{x'}$ (with $\cl(Y')=1\in \h^0(X,\dZ/n)$), and
$(\text{id},u)$, which falls under (B).

\paragraph{Proof of (iii)}
Let $\delta:X\to X\times_S X$ be the diagonal inclusion. One uses the formula
$\delta^\ast(Y_1\times_S Y_2) = Y_1\cap Y_2$ to reduce to
(B) and (C).

\paragraph{Proof of (i)}
Base change $S_\text{red}\to S$ replaces $X$ by $X_\text{red}$; this
reduces us to assuming $X$ and $S$ reduced. A localization near
the generic points of $Y$ then reduces us to the case where
$Y_\text{red}$ is irreducible and the complete intersection of $c$
divisors $D_i$. The same localization, and the definition of $\tr$ in the
quasi-finite case, then shows that
$\cl(Y,\sK)=\operatorname{lg}\cl(Y_\text{red},\sO_{Y_\text{red}})$, where
$\operatorname{lg}$ is the length of $\sK$ at the generic point of $Y$.
By (iv), one therefore has
\[
  \cl(Y,\sK) = \operatorname{lg} \prod_i \cl(D_i,\sO_{D_i})
\]
and one concludes by \ref{IV:2-3-6}.




\subsubsection{Remark}\label{IV:2-3-9}

For $S$ the spectrum of a field, and the class of a cycle being defined as
in \ref{IV:2-2-10}, assertion (iii) says that, if two cycles $Y'$ and $Y''$ meet
in the right dimension, then
\[
  \cl(Y'\cap Y'') = \cl(Y')\smallsmile \cl(Y'')\text{,}
\]
provided the multiplicities of the components of $Y'\cap Y''$ are
computed as alternating sums of $\tor$.




\subsubsection{Remark}\label{IV:2-3-10}

Let $X/\spec(k)$, $k$ algebraically closed. If two cycles of codimension $c$
in $X$ are algebraically equivalent, it follows from the existence of the
relative theory above that the images in $\h^{2 c}(X,\dZ/n(c))$ of
their classes are the fibers, at two points of a suitable connected $S$, of a
section over $S$ of the constant sheaf $\h^{2 c}(X,\dZ/n(c))$ (the sheaf
$\eR^{2 c} f_\ast \dZ/n(c)$ for $f=\operatorname{pr}_2:X\times S\to S$): two
algebraically equivalent cycles have the same class in
$\h^{2 c}(X,\dZ/n(c))$.




















\section{Application: the Lefschetz trace formula in the proper smooth case}\label{IV:3}




\subsection{}\label{IV:3-1}

Let $X$ and $Y$ be proper smooth algebraic varieties over an algebraically
closed field $k$, purely of dimensions $N$ and $M$. Via
the K\"unneth isomorphism
$\h^\bullet(X\times Y)=\h^\bullet(X)\otimes \h^\bullet(Y)$, the trace
morphism $\h^\bullet(X\times Y)(M)\to \h^\bullet(X)$ is none other than
$\h^\bullet(X)\otimes ($the trace morphism $\h^\bullet(Y)(M)\to \dQ_\ell)$. The
morphism being homogeneous of even degree, there is no sign
problem. We denote it by $\int_Y$.




\subsection{}\label{IV:3-2}

A class $\eta\in \h^{2 q}(X\times Y)(q)$ defines a morphism of degree
$2(q-N)$ $\eta_{XY}:\h^\bullet(Y)(N-q)\to \h^\bullet(X)$, by the formula
\[
  \beta \mapsto \int_Y \eta\cdot \operatorname{pr}_2^\ast \beta \text{.}
\]




\begin{proposition_}\label{IV:3-3}
Let $p$ and $q$ be two integers, with $p+q=N+M$, and
$\varepsilon\in \h^{2 p}(X\times Y)(p)$, $\eta\in \h^{2 q}(X\times Y)(p)$. The
trace of $\eta_{X Y}\varepsilon_{Y X}:\h^\bullet(X)\to \h^\bullet(X)$ being
understood in the sense of \ref{IV:1-3}, one then has
\[
  \tr_{X\times Y}(\eta\cdot \varepsilon) = \tr\left(\eta_{X Y}\varepsilon_{Y X},\h^\bullet(X)\right) \text{.}
\]
\end{proposition_}

To exorcise signs, it is convenient to retain from the grading
of $\h^\bullet$ only the underlying $\dZ/2$-grading. To avoid carrying
Tate twists along, one further fixes an isomorphism
$\dQ_\ell(1)=\dQ_\ell$. Let $\alpha_Y$ be the homomorphism
$\h^\bullet(Y)\to \h^\bullet(Y)^\vee$ (an isomorphism) making the diagram
\[\xymatrix{
  \h^\bullet(Y)\otimes \h^\bullet(Y) \ar[r] \ar[d]^-{\alpha_Y\otimes 1}
    & \h^\bullet(Y) \ar[r]^-{\tr}
    & \dQ_\ell \ar@{=}[d] \\
  \h^\bullet(Y)^\vee \otimes \h^\bullet(Y) \ar[rr]^-{\text{(\ref{IV:eq:1-3-1})}}
    & & \dQ_\ell \text{.}
}\]
commute. The map $\eta_{X Y}$ is the image of $\eta$ under
\[\xymatrix{
  \h^\bullet(X\times Y) = \h^\bullet(X)\otimes \h^\bullet(Y) \ar[r]^-{\alpha_Y}
    & \h^\bullet(X) \otimes \h^\bullet(Y)^\vee \ar@{=}[r]^-{\text{(\ref{IV:eq:1-3-2})}}
    & \hom\left(\h^\bullet(X),\h^\bullet(Y)\right) \text{,}
}\]
and likewise for $\varepsilon$. Definition \ref{IV:eq:1-3-4} of the trace
then reduces \ref{IV:3-3} to the formula
$\tr_{X\times Y}=\tr_X\otimes \tr_Y$.




\subsection{Remark}\label{IV:3-4}

The proof of \ref{IV:3-4} does not use that $\alpha_X$ and $\alpha_Y$ are
isomorphisms. It holds without assuming $X$ and $Y$ smooth. However, outside
the smooth case, it is hard to find \emph{cohomology}
classes to which one would want to apply \ref{IV:3-3}.




\subsection{Remark}\label{IV:3-5}

If $\varepsilon,\eta$ are the cohomology classes of algebraic cycles
$A$ and $B$ on $X\times Y$, of codimension $p$ and $q$, and $A\cap B$ is of
dimension $0$, then $\tr_{X\times Y}(\eta\varepsilon)$ is the intersection
multiplicity $A\cdot B$, computed by an alternating sum of $\tor$
(\ref{IV:2-3-9}).




\begin{proposition_}\label{IV:3-6}
For $q=M$ and $\eta=\cl(B)$ the class of a subscheme $B\subset X\times Y$,
finite over $X$, the map $\eta_{X Y}$ is the composite
\[\xymatrix{
  \eta_{X Y} : \h^\bullet(Y) \ar[r]^-{\operatorname{pr}_2^\ast}
    & \h^\bullet(B) \ar[r]^-{\tr_{B/X}}
    & \h^\bullet(X) \text{.}
}\]
\end{proposition_}

By \ref{IV:3-2} and \eqref{IV:eq:2-3-3-1} for $\sK=\sO_B$, one has
indeed
\[
  \eta_{X Y}(\beta) = \tr_{X\times Y/X}\left(\cl(B)\cdot \operatorname{pr}_2^\ast\beta\right) = \tr_{B/X}\left(\operatorname{pr}_2^\ast \beta\right) \text{.}
\]

The most important case is that where $B$ is the graph of a map
$f:X\to Y$. In this case, \ref{IV:3-6} shows that $\eta_{X Y}=f^\ast$.




\begin{corollary_}\label{IV:3-7}
Let $f$ be an endomorphism of $X$ proper and smooth over an algebraically
closed field $k$. Assume the fixed points of $f$ isolated. Then the
trace $\sum (-1)^i \tr(f^\ast,\h^i(X))$ is the number of fixed points of $f$,
each counted with its multiplicity.
\end{corollary_}

In \ref{IV:3-5}, one takes $A=$graph of the identity, $B=$graph of $f$ and one
applies \ref{IV:3-3}, \ref{IV:3-6}.




\begin{corollary_}\label{IV:3-8}
Let $X$ be a curve over $k$ algebraically closed, deduced by scalar
extension from $X_0$ over $\dF_q$, and $f$ the Frobenius morphism. Then,
$\sum (-1)^i \tr\left(f^\ast,\h_c^\bullet(X)\right)$ is the number of fixed
points of $f$.
\end{corollary_}

One may replace $X$ by $X_\text{red}$, hence assume $X$ reduced. Let $U$
be the open subset of $X$ where $X$ is smooth of dimension $1$, and $\bar U$ the
smooth projective completion of $U$. The long exact cohomology
sequences for $(X,U)$ and $(\bar U,U)$ give
\[
  \tr\left(f^\ast,\h_c^\bullet(X)\right) = \tr\left(f^\ast,\h_c^\bullet(U)\right) + \tr\left(f^\ast,\h_c^\bullet(X\setminus U)\right)
\]
and
\[
  \tr\left(f^\ast,\h_c^\bullet(\bar U)\right) = \tr\left(f^\ast,\h_c^\bullet(U)\right) + \tr\left(f^\ast,\h_c^\bullet(\bar U\setminus U)\right) \text{.}
\]
The same formulas hold for the numbers of fixed points. Formula
\ref{IV:3-8} being clear for a scheme of dimension $0$, this reduces
proving \ref{IV:3-8} to $\bar U$. This case falls under \ref{IV:3-7}.
The fact that $d f=0$ ensures that the fixed points are all of
multiplicity one.




\subsection{Remark}\label{IV:3-9}

For $X$ of dimension $\leqslant 1$ over $k$ algebraically closed, and
$f:X\to X$ with $d f=0$, it is not hard to check that
\ref{IV:3-8} still holds.
